\section{Inputs: the characteristics and how they are prepared}
\label{sec:features}

\dec{2026-10-05, Q26} 40 characteristics and 57 model inputs, in four blocks. Implemented as
\texttt{feature\_set="q26"}; the old 14 price features survive only as \texttt{"legacy14"}.

\subsection{The selection rule}

\textbf{Exactly two representatives per theme} of Jensen, Kelly \& Pedersen (2023, ``JKP''), who cluster 153
published characteristics into 13 themes. Representatives were chosen to be robust in large caps, well
documented, computable from CRSP and Compustat, and not near-copies of each other. PCA on training data is
used only as a redundancy diagnostic.

\begin{whybox}
\begin{itemize}
\item \textbf{A rule, fixed before any result, instead of a search.} Picking features by how well they
  predict would be selection on outcomes. A rule stated in advance cannot be tuned to the data.
\item \textbf{Fewer, well-chosen inputs suit a weak signal and 500 stocks.} Hundreds of characteristics raise
  estimation variance and give the experts more room to relearn general signal (the identification risk of
  Q10). Two per theme covers every documented source of return predictability with about 40 inputs.
\item \textbf{Two, not one, per theme}: one representative can be noisy or fail in a subperiod; two let the
  network combine them, while the fixed count keeps themes balanced so no theme dominates by sheer number
  of inputs.
\item \textbf{Why not PCA as the selector}: principal components mix themes into uninterpretable
  combinations, change from fold to fold, and would have to be re-estimated inside each training block. As a
  diagnostic it still shows whether two representatives are redundant.
\end{itemize}
\end{whybox}

\subsection{Block A: JKP themes (26 characteristics)}

{\small
\begin{tabularx}{\textwidth}{@{}l X X l@{}}
\toprule
Theme & Representative 1 & Representative 2 & Data\\
\midrule
Short-term reversal & 5-day return & 20-day return & CRSP\\
Momentum & 12-month return skipping the last month & price / 52-week high & CRSP\\
Low risk & 21-day volatility & beta (Frazzini--Pedersen) & CRSP\\
Size & log market cap & Amihud illiquidity, 126 days & CRSP\\
Seasonality & same-calendar-month return, years 2--5 back & coskewness with the market, 21 days & CRSP\\
Value & book-to-market & earnings-to-price & both\\
Profitability & quarterly ROE & operating cash flow / assets & Compustat\\
Quality & gross profitability / assets & earnings consistency & Compustat\\
Investment & asset growth & sales growth & Compustat\\
Accruals & operating accruals & total accruals & Compustat\\
Debt issuance & 3-year debt growth & net operating assets / assets & Compustat\\
Low leverage & net debt / market cap & cash / assets & both\\
Profit growth & standardised unexpected earnings (SUE) & revenue surprise & Compustat\\
\bottomrule
\end{tabularx}}

\subsection{Block B: short-horizon market themes (12 characteristics)}

{\small
\begin{tabularx}{\textwidth}{@{}l X X@{}}
\toprule
Theme & Representative 1 & Representative 2\\
\midrule
Volatility dynamics & idiosyncratic volatility, 21 days & volatility shock: 5-day / 60-day volatility\\
Tails and asymmetry & realised skewness, 21 days & maximum daily return, 21 days\\
Volume and liquidity & abnormal volume (today vs trailing average) & quoted spread, 21-day average\\
Price path & overnight (close-to-open) return, 20 days & price / 50-day moving average\\
Earnings timing & announcement expected in the next 5 days (0/1) & trading days since last announcement\\
Return dynamics & variance ratio, 5-day vs 1-day, 60 days & return--volume correlation, 60 days\\
\bottomrule
\end{tabularx}}

\begin{whybox}[Why a short-horizon block]
The JKP themes come from the \emph{monthly} anomaly literature, and 8 of 13 are accounting-based.
Fundamentals barely change from week to week, while the signals known to predict weekly returns are
market-based: price trends, liquidity and volatility dominate the feature importance of Gu, Kelly \& Xiu.
Applying the same two-per-theme rule to six short-horizon themes, each with evidence at a daily or weekly
horizon, matches the inputs to the target's horizon without abandoning the rule. Several were also
confirmed in Tom's own Trexquant work (volatility term structure, the overnight/intraday split, days to
earnings, volume relative to its own average). Levels are expressed as surprises against the stock's own
history (volume, volatility) because a raw level mostly acts as a static tilt.
\end{whybox}

\subsection{Block C: industry (17 inputs) and Block D: flags (2)}

\begin{itemize}
\item 11 \textbf{GICS sector dummies} (0/1); \dec{2026-10-05} GICS from Compustat.
\item \textbf{Industry momentum} (the sector's 12-month return skipping the last month; Moskowitz \&
  Grinblatt 1999) and \textbf{within-industry reversal} (the stock's 20-day return minus its sector's).
\item 4 \textbf{within-sector ranks}: book-to-market, earnings-to-price, quarterly ROE, operating cash
  flow / assets.
\item 2 \textbf{missing-data flags}: one for price/volume inputs, one for fundamentals.
\end{itemize}
Total: $26+12+11+2+4+2=57$ inputs.

\begin{whybox}
\begin{itemize}
\item \textbf{Industries differ in what ``cheap'' or ``profitable'' means.} A bank's book-to-market is not
  comparable with a software firm's. A universe-wide rank of such ratios partly just sorts industries;
  ranking within the sector compares a firm with its peers.
\item \textbf{Sector dummies} let the experts learn sector-specific shifts, which matters more now that the gate
  has an industry level.
\item \textbf{Industry momentum} is a documented, separate source of predictability. \textbf{Within-industry
  reversal} separates a stock's own overreaction from news that moved its whole industry, which is the
  part of short-term reversal most associated with liquidity provision.
\item \textbf{GICS, not ICB}: CRSP stops filling ICB for S\&P 500 members in October 2023, and Compustat's
  GICS history is dated, so each stock gets the sector it had \emph{at the time}, not today's.
\end{itemize}
\end{whybox}

\subsection{Scaling: ranks mapped to $[-1,1]$}

\dec{2026-10-05, Q26 part 2} On each date each characteristic is ranked across the universe (ties get the
average rank) and mapped linearly:
\begin{equation}
x_{i,t} = \frac{2\,(\mathrm{rank}_{i,t}-1)}{N_t-1}-1 \in [-1,1].
\end{equation}
Dummies and flags are not ranked.

\begin{whybox}
\begin{itemize}
\item \textbf{Robust to outliers and scale.} Characteristics like illiquidity or accruals have extreme tails
  that destabilise neural network training; ranks remove them.
\item \textbf{Comparable over time.} A book-to-market of 0.5 means different things in 2000 and 2021; being in
  the top decile means the same thing.
\item \textbf{Only cross-sectional information survives}, which is exactly what the market-neutral target
  measures. This is the convention of Gu, Kelly \& Xiu.
\item \textbf{$[-1,1]$ rather than $[-0.5,0.5]$}: input variance $1/3$ instead of $1/12$, closer to the unit
  variance that standard weight initialisations assume, so the first layer starts in a sensible range.
\end{itemize}
\giveup{the information in absolute levels and in how far apart stocks are, which ranking discards.}
\end{whybox}

\subsection{Missing values: causal handling}

\dec{2026-10-05, Q26 part 3}
\begin{enumerate}
\item Price/volume windows need at least 80\% of their days present.
\item Fundamentals carry the firm's last point-in-time value forward for at most 12 months.
\item Anything still missing is set to 0 after ranking (the median) and the family's flag is set to 1.
\item Rows are no longer dropped for a missing characteristic.
\end{enumerate}

\begin{whybox}
\begin{itemize}
\item \textbf{Never use the future.} Interpolation fills a gap with a later value, which is look-ahead.
  Every rule here uses only data available at $t$.
\item \textbf{Keep every stock.} At forecast time every stock in the universe needs a forecast. A training rule
  that drops incomplete rows trains on a different population from the one it is used on, and missingness
  is not random (new listings, distressed firms).
\item \textbf{Carrying fundamentals forward is barely an imputation}: between two reports the latest filing is
  exactly what the market knows. The 12-month cap stops a stale value standing in for a firm that stopped
  reporting.
\item \textbf{80\% windows}: an average over 16 of 20 days estimates the same quantity with little extra noise;
  requiring all 20 would discard many stock-days around trading halts.
\item \textbf{Zero plus a flag}: in rank space zero is the median, the neutral ``no information'' value, and
  the flag lets the model learn a separate shift for stocks with missing data instead of the fill
  distorting the feature.
\end{itemize}
Four Compustat items (IVAOQ, IVSTQ, MIBQ, PSTKQ) count as 0 when blank in total accruals and net operating
assets: blank there usually means ``none'', and the strict rule left those features on only about 5\% of
rows (79\% after the change).
\end{whybox}

\subsection{Point-in-time fundamentals}

\dec{2026-10-05, Q26 part 5} A quarter's numbers become usable \textbf{one trading day after the later of its
report date and its filing date}; surprises (SUE, revenue surprise) are dated by the report date; flows are
trailing four-quarter sums; market values come from CRSP on day $t$.

\begin{whybox}
\begin{itemize}
\item \textbf{The quarter end is not when investors learn the numbers.} Firms report weeks later, and some items
  appear only in the filing. Using the quarter end plus a fixed lag either leaks information (lag too short)
  or wastes it (too long). Report and filing dates say exactly when the numbers became public.
\item \textbf{One trading day after}, because many announcements come after the close; this removes any
  ambiguity about whether the close of the announcement day could have reacted.
\item \textbf{Surprises by the report date}, because the announcement is the event the market reacts to.
\item \textbf{Trailing four-quarter sums} remove seasonality from flows such as earnings and cash flow.
\end{itemize}
\giveup{restatements: standard Compustat values may be revised after first release (a known limitation).}
\end{whybox}

Every characteristic has a no-look-ahead test: its value at $t$ computed from data truncated at $t$ equals
its value from the full data.

\subsection{What is deliberately not an input}

\begin{itemize}
\item \textbf{Market-wide series} (VIX, rates, FOMC dates, macro data, the market return): identical across
  stocks on a date, so ranking removes them. They belong in the gate. This is also what the closest
  precedent found on its volatility target: appending regime variables to the input made 30 of 30 seeds
  collapse, while feeding them through the gate made none collapse.
\item \textbf{Intraday data, trade counts} (28\% coverage), analyst estimates, options, news: not in the data.
\item Near-copies of included features (position in the 20-day range, high/low ratio, up/down volatility
  ratio).
\end{itemize}

\begin{keybox}{Scoping future alphas}
A new input must be one number per stock per day, known by the close of $t$, varying across stocks; only
its ordering matters. Signals that pay off over 1--10 days fit best. A machine-learning alpha must be
trained inside the same walk-forward folds on out-of-sample predictions only (Improvements item E2).
Inputs added after seeing test results are a forking-paths problem.
\end{keybox}
